% © Copyright 2012 Joachim Breitner
% Distributed under the terms of the Creative Commons Attribution license.
\documentclass{beamer}
%\documentclass{article}
%\usepackage{beamerarticle}

\newlength{\px}
\setlength{\px}{0.0009765625\paperwidth}

\usetheme[footline=authortitle,compress]{Nomeata}

\setbeamertemplate{headline}{}

\AtBeginSection[]
{
%  \begin{frame}<beamer>
%    \tableofcontents[currentsection,currentsubsection]
%  \end{frame}
}

%\usecolortheme{orchid}
\usecolortheme{crane}

\usepackage[T3,OT1]{fontenc}
\usepackage[ngerman]{babel}
\usepackage[utf8]{inputenc}
\usepackage{tikz}
\usetikzlibrary{decorations.pathreplacing,calc}
\usetikzlibrary{positioning,calc,decorations,decorations.pathmorphing,shapes.geometric,matrix}
\usetikzlibrary{shapes.geometric}
\usetikzlibrary{calc}
\usetikzlibrary{shapes.multipart}
\usetikzlibrary{arrows}
\usepackage{hyperref}
\usepackage{movie15}

\usepackage{mathtools}
\usepackage{amssymb}

\usepackage{listings}

\definecolor{light-gray}{gray}{0.95}
\lstdefinestyle{haskell}{
        language=Haskell
        ,literate=
                {ö}{{\"o}}1
                {ä}{{\"a}}1
                {ü}{{\"u}}1
        ,columns=flexible
        ,basewidth={.365em}
        ,keepspaces=True
        ,texcl=true
        ,basicstyle=\small\sffamily
        ,stringstyle=\itshape
        ,showstringspaces=false
        ,keywords={module,where,open,import,using,renaming,to,data,let,in,with}
        ,belowskip=0pt
        ,backgroundcolor=\color{light-gray}
        ,frame=single
        ,xleftmargin=2em
        ,xrightmargin=2em
        }


% \pause mit verstecken
\newcommand{\hide}{\onslide+<+(1)->}

\title{Haskell Bytes}
\subtitle{Eine geführte Tour durch den Hauptspeicher eines Haskell-Programms}
\author{Joachim Breitner}
%\institute{Karlsruher Institut für Technologie, IPD Snelting}
\titlegraphic{\includegraphics[width=11em]{HaskellBytesLogo}}
\date{MRMCD 12, 8. September 2012}

\DeclareTextSymbol\textlambda{T3}{171}           % Lambda
\DeclareTextSymbolDefault\textlambda{T3}
\DeclareUnicodeCharacter{03BB}{\textlambda}


\begin{document}

\frame[plain]{\titlepage \vfill
\raggedleft
\includegraphics[width=5em]{Telekom}

}
\only<article>{\maketitle}

\only<presentation>{
\centering
\frame[plain]{\textit{Demonstration}}
}

\begin{frame}
\frametitle{Drei Arten von Closures}
\begin{center}
\begin{tikzpicture}[
        closure/.style={shape=rectangle, rounded corners, minimum width=2.8cm,fill=blue!40,font=\bfseries},
        prop/.style={},
	brace/.style={decoration={brace,amplitude=0.5em},decorate,thick}
	]
\hide
\node[closure] (con) {Konstruktor};
\hide
\node[closure] at (3,0) (fun) {Funktion};
\hide
\node[closure] at (6,0) (thunk) {Thunk};
\hide
\draw [brace] (con.south east) -- (con.south west);
\path (con) ++(0,-0.7) node[prop] {Parameter};
\hide
\path (con) ++(0,-1.2) node[prop] {Welcher Konstr.};
\hide
\draw [brace] (thunk.south east) -- (fun.south west);
\path (4.5,0)  ++(0,-0.7) node[prop] {Code};
\hide
\path (4.5,0)  ++(0,-1.2) node[prop] {Freie Variablen};
\hide
\draw [brace] ($(thunk.south east)+(0,-1.5)$) -- ($(con.south west)+(0,-1.5)$);
\path (3,-1.5)  ++(0,-0.7) node[prop] {Größe};
\hide
\path (3,-1.5)  ++(0,-0.7) ++(0,-0.5) node[prop] {Layout};
\hide
\path (3,-1.5)  ++(0,-0.7) ++(0,-0.5) ++(0,-0.5) node[prop] (typ) {Typ};
\hide
\draw [line cap=rounded, ultra thick, color=blue] (typ.south west) -- (typ.north east);
\onslide+<0->
\end{tikzpicture}
\end{center}
\end{frame}

\only<presentation>{
\centering
\frame[plain]{\textit{Investigation}}
}

\only<presentation>{
\centering
\frame[plain]{
\includemovie[autoplay,mouse=true,mimetype=video/webm,poster=GC-frame1.png,inline=false,controls=true]{8cm}{6cm}{GC.webm}
}
}

\only<presentation>{
\centering
\frame[plain]{\textit{Investigation}}
}


\def\ux{2.7cm}\def\uy{0.6cm}
\begin{frame}
\frametitle{Der Info-Pointer}
\begin{center}
\begin{tikzpicture}[x=\ux, y=\uy,
        word/.style={shape=rectangle, draw, minimum width=\ux, minimum height=\uy,font=\small}
        ,halfword/.style={shape=rectangle, draw, minimum width=0.5*\ux, minimum height=\uy,font=\small}
        ,>=latex]
\draw (0,0) rectangle +(1,1) node[midway] (ip) {Info-Pointer}
    ++(1,0) rectangle +(1,1) node[midway] {Pointer}
%    ++(1,0) rectangle +(1,1) node[midway] {Pointer}
    ++(1,0) +(0.1,0.4) node {$\cdots$}
    ++(0.2,0) rectangle +(1,1) node[midway] {Nicht-Pointer}
 %   ++(1,0) rectangle +(1,1) node[midway] {Nicht-Pointer}
    ++(1,0) +(0.1,0.4) node {$\cdots$};

\begin{scope}[yshift=-0.7cm, xshift=1.2*\ux]
\draw
  (0,0) node[word] (tbl) {Code-Pointer}
++(-0.25,-1) node[halfword] {\#ptr} ++(+0.5,0) node[halfword] {\#nptr} ++(-0.25,0)
++(-0.25,-1) node[halfword] {Typ} ++(+0.5,0) node[halfword] {SRT} ++ (-0.25,0)
++(0,-1) node[word] {ggf. Konstruktor}
++(0,-1) node[word] {ggf. Stelligkeit};
\end{scope}
\draw[*->] (ip.south) |- (tbl.west);
\draw[*->] (tbl.east) -- ++(1.5cm,0) node[right] (code) {Entry code};
\draw (code) ++(-0.5,-1) -- ++(0,1.5) -- ++(1,0) -- ++(0,-1.5);
\end{tikzpicture}
\end{center}
\end{frame}

\begin{frame}
\frametitle{Der echte Info-Pointer}
\begin{center}
\begin{tikzpicture}[x=\ux, y=\uy,
        word/.style={shape=rectangle, draw, minimum width=\ux, minimum height=\uy,font=\small}
        ,halfword/.style={shape=rectangle, draw, minimum width=0.5*\ux, minimum height=\uy,font=\small}
        ,>=latex]
\draw (0,0) rectangle +(1,1) node[midway] (ip) {Info-Pointer}
    ++(1,0) rectangle +(1,1) node[midway] {Pointer}
%    ++(1,0) rectangle +(1,1) node[midway] {Pointer}
    ++(1,0) +(0.1,0.4) node {$\cdots$}
    ++(0.2,0) rectangle +(1,1) node[midway] {Nicht-Pointer}
%    ++(1,0) rectangle +(1,1) node[midway] {Nicht-Pointer}
    ++(1,0) +(0.1,0.4) node {$\cdots$};

\begin{scope}[yshift=-0.7cm, xshift=1.2*\ux]
\draw
  (0,0) node[word] {ggf. Konstruktor}
++(0,-1) node[word] {ggf. Stelligkeit}
++(-0.25,-1) node[halfword] {\#ptr} ++(+0.5,0) node[halfword] {\#nptr} ++(-0.25,0)
++(-0.25,-1) node[halfword] {Typ} ++(+0.5,0) node[halfword] {SRT} ++ (-0.25,0)
++(0,-1) node[word,draw=none] (code) {Entry code};
\end{scope}
\draw[*->] (ip.south) |- (code.west);
\draw (code) ++(-0.5,-1) -- ++(0,1.5) -- ++(1,0) -- ++(0,-1.5);
\end{tikzpicture}
\end{center}
\end{frame}

\setbeamercolor{normal text}{bg=black}
\setbeamertemplate{background canvas}{}
\frame<presentation>[plain]{
\strut
\vfill
\small{© 2012 Joachim Breitner.\\
Distributed under the terms of the Creative Commons Attribution license.}
}
\end{document}

